% ----------------------------------------------------------
\chapter{Tabela de \textit{MIBs} SNMP Ruckus Wireless}
\label{an:oids_ruckus}
% ----------------------------------------------------------

Este anexo apresenta os identificadores de objetos (OIDs) da \textit{MIB} (\textit{Management Information Base}) proprietária da Ruckus Wireless consultados periodicamente via SNMP pelo coletor desenvolvido neste trabalho. Todas as métricas pertencem à tabela de clientes da controladora, cujo OID base é \texttt{1.3.6.1.4.1.25053.1.2.2.1.1.3.1.1}. Cada métrica corresponde a uma coluna dessa tabela, identificada pelo número acrescentado ao OID base, e o endereço MAC do cliente compõe o índice de cada linha. O ruído de fundo não é consultado diretamente: ele é derivado como a diferença entre o sinal recebido e a relação sinal-ruído (SNR).

\begin{table}[htb]
\IBGEtab{%
  \caption{Mapeamento de OIDs SNMP Utilizados na Coleta Ruckus}%
  \label{tab:oids_snmp_ruckus}%
}{%
  \resizebox{\textwidth}{!}{%
    \begin{tabular}{lll}
      \toprule
      Métrica & OID SNMP consultado & Descrição \\
      \midrule
      Endereço MAC do AP & \texttt{<base>.2} & AP ao qual o cliente está associado \\
      Banda do rádio & \texttt{<base>.6} & Identificador da banda de operação (2,4 GHz ou 5 GHz) \\
      Canal de RF & \texttt{<base>.7} & Canal do rádio ao qual o cliente está associado \\
      Pacotes transmitidos & \texttt{<base>.10} & Contador cumulativo de pacotes transmitidos (TX) \\
      Bytes transmitidos & \texttt{<base>.11} & Contador cumulativo de bytes transmitidos (TX) \\
      Bytes recebidos & \texttt{<base>.13} & Contador cumulativo de bytes recebidos (RX) \\
      Retransmissões TX & \texttt{<base>.14} & Contador cumulativo de quadros retransmitidos \\
      Tempo de associação & \texttt{<base>.15} & Tempo de conexão do cliente (\textit{uptime}) \\
      Taxa de ligação & \texttt{<base>.17} & Velocidade física negociada com o cliente (Kbps) \\
      SNR & \texttt{<base>.21} & Relação sinal-ruído do cliente (dB) \\
      Sinal recebido (RSSI) & \texttt{<base>.81} & Intensidade do sinal recebido do cliente (dBm) \\
      \bottomrule
    \end{tabular}%
  }%
}{%
  \fonte{Produzido pelos autores a partir do script de coleta. \texttt{<base>} = \texttt{1.3.6.1.4.1.25053.1.2.2.1.1.3.1.1}.}%
}
\end{table}
